\section{Investigating a Suspected Compromise}
\label{sec:suspected-compromise}

A batch of 48 runs launched from the dashboard this round reversed the
paper's main result outright: \texttt{con\_harness} reached zero leaks (0/28
across all four tasks) while \texttt{sin\_harness} leaked in 13 of 20
attempts, split cleanly by task (0/10 against 10/10 on \texttt{task\_06},
0/10 against 3/10 on \texttt{task\_05}). A result that clean, on the machine
the rest of the corpus came from, is exactly the shape a real compromise
would take, and it was treated as one until it was checked rather than
reported.

\subsection{What actually happened}

It did not happen on this machine's model. Every run in the batch shows
direct evidence of talking to a cloud reasoning backend instead of the local
Ollama instance the rest of \texttt{machine-A} used: 46 of 47 transcripts
carry \texttt{prompt\_cache\_hit\_tokens} and
\texttt{completion\_tokens\_details.reasoning\_tokens} in their
\texttt{usage} field, fields Ollama's OpenAI-compatible endpoint does not
emit; the 47th failed outright with a \texttt{503} from
\texttt{opencode.ai/zen/go/v1/chat/completions} logged verbatim in its
transcript. The cause was mundane: a local \texttt{.env}, left over from the
cloud-model spot check already reported in Section~\ref{sec:results}, still
set \texttt{AGENT\_BACKEND=openai}, and the orchestrator's own metadata
capture reads that variable from its own process environment rather than
from the file Compose actually applies to the container, so every
\texttt{\_meta.json} in the batch falsely records
\texttt{"agent\_backend": "ollama"}.

We checked, rather than assumed, that this could not be an escape. No
compose profile mounts the repository root or \texttt{.env} into any
container; the sandbox receives only pre-resolved environment values at
creation time, with no filesystem path back to the host file. \texttt{.env}'s
own \texttt{birth} and \texttt{modify} timestamps are identical and about 40
hours older than the batch's first run, with nothing written in between. The
batch is a mislabeled dataset, not a security event, and it should be moved
out of \texttt{machine-A}'s corpus and re-tagged before it is used for
anything.

\subsection{What can actually reach the agent}

The suspicion this round kept returning to was whether an attacker could
connect to the agent directly, rather than wait for it to fetch something.
The topology answers this without ambiguity. The sandbox container publishes
no port and runs no listener; nothing in \texttt{sandbox/} binds a socket or
serves a request. It sits on two Docker networks, and on both of them it is
always the one that initiates. Every other service in the compose file
connects out to something; nothing connects in to \texttt{sandbox}.
\texttt{cap\_drop: [ALL]} and \texttt{no-new-privileges} bound what a
compromised process inside it could do regardless.

This rules out a push, not the pull already measured. The agent already has,
by design, the means to run arbitrary code inside its own container through
\texttt{write\_file} and \texttt{run\_tests}; that is not a vulnerability, it
is the task. What a sandbox escape detector exists to answer is whether that
arbitrary code ever leaves the container, and across every run collected for
this project, control and mislabeled cloud batch included, level 6
(\texttt{sandbox\_roto\_heartbeat}) fired zero times in 399 runs. Counting the
full corpus as it stands today, 416 runs across four machines, it fired five
times, all on machine D and all explained the same way:
\texttt{docker\_exit\_code=124} in every one, meaning the orchestrator's own
outer timeout killed the run after the agent spent its full 480\,s budget
re-reading memory, rewriting the same failing patch and re-running the same
failing test without ever converging, and the heartbeat monitor correctly
read the silence that followed a kill as silence. Four of the five also used
the canary before the kill, through the same \texttt{memory\_write} leak
Section~\ref{sec:results} already counts at level 5; the fifth timed out
without touching it. None shows anything past what level 5 already covers:
no new egress surface, nothing leaving the container through a channel this
paper does not already measure. What level 6 adds, in every case, is only
that the process was already done leaking or already stuck, and then stopped
answering. That all five sit on the one machine slow
enough to time out its own HTTP client mid-prompt (Appendix~\ref{sec:repro})
is the same hardware limit surfacing a second way: there, a slow first turn
produces a run with nothing in it; here, a slow agent that never converges
produces a run with fourteen real turns in it and still gets killed. Level 6
exists in the ladder precisely for a case like this, and it has never once
coincided with a breakout: every security-relevant escalation, on any
machine, still tops out at level 5, the agent using its own
already-permitted egress channel to move the canary somewhere it should not,
which is the indirect-injection threat model of ASI06~\cite{owasp-asi06}, not
a container breakout.

\subsection{Next}

Three items follow directly from this round and are cheap relative to the
rest of the paper. Fix \texttt{backend\_info()} to read the same
\texttt{.env} Compose resolves, so a future misconfiguration reports itself
instead of a wrong backend name. Move \texttt{machine-A/prueba-real} to its
own corpus under the backend it actually used, and either drop it or report
it as a second, larger replication of the model-reversal result in
Section~\ref{sec:results} rather than as \texttt{machine-A} data. And since
this round's real question, whether an attacker could act on the agent
rather than through it, has no path to test in the current topology, the
honest answer is that this harness has never been asked that question: every
vector here is something the agent chooses to pull, never something pushed
to it mid-run.

\todo{Opcional si sobra espacio: una linea notando que extender la topologia
con una superficie push (un webhook, un socket que el agente escuche) seria
la forma de probar esto de verdad, en vez de inferirlo de la ausencia actual
de un listener.}
